\chapter*{Hamiltonian jets and boundary quantization}
\addcontentsline{toc}{chapter}{Hamiltonian jets and boundary quantization}

Quantization of a symplectic boundary changes the product of observables while retaining the classical Poisson bracket to first order. The formal symplectic plane gives a concrete setting in which to compare Hamiltonian operations with operations on a quantized brane. Pointed Hamiltonians vanish to order at least two; translations belong to a separately adjoined affine sector. Finite Hamiltonian jets and their shifted cotangent fields provide the local algebraic models.

The universal boundary brane leads to a Rees enveloping algebra. Its relative Hochschild complex is compared with a Chevalley complex whose differential carries the same Rees parameter. Duflo quantization gives the separate invariant comparison, and the quadratic Moyal calculation determines the oscillator Casimir quotient. These finite constructions distinguish the underlying cochain comparison from higher operations on it.

The all-arity open--closed problem asks for compatible closed and mixed operations on the boundary algebra. Its stated refinement supplies the operadic maps and coherence homotopies. Passing through the Hamiltonian jet limit also requires actual cochain transition systems and their continuous limit identifications. A geometric brane realization must supply the boundary module and compare its centre action with the chosen observable or trace maps; global realization further requires descent and quantization data.
